\documentclass{article}
\usepackage{amsmath}
\usepackage{amssymb}
\usepackage{amsthm}
\usepackage{enumitem}
\usepackage{titling}
\usepackage{graphicx}
\usepackage{tikz-cd}
\usepackage[hidelinks]{hyperref}
\usepackage{nameref}
\usepackage{fancyhdr}
\usepackage[a4paper, margin=1in]{geometry}
\usepackage{refcount}
\usepackage{physics}

% \tikzcdset{scale cd/.style={every label/.append style={scale=#1},
%     cells={nodes={scale=#1}}}}

\newtheorem{theorem}{Teorema}
\newtheorem{corollary}{Corollario}[theorem]
\newtheorem{osservazione}{Osservazione}
\newtheorem{lemma}{Lemma}
\newtheorem{proposition}{Proposizione}
\newtheorem{definition}{Definizione}
\newtheorem{axiom}{Assioma}

\setlength{\parindent}{0pt}

\begin{document}
\section{Numeri complessi}
    I numeri complessi $\mathbb{C} = \mathbb{R}(i)$ sono un'estensione del campo $\mathbb{R}$ secondo l'aggiunta di un nuovo elemento, $i$, chiamato unità immaginaria.
    Il campo che ne risulta, $\mathbb{C}$, è un campo \textit{algebricamente chiuso}.
    Possiamo definire l'unità immaginaria come 
    \[
        i^{2} = -1
    \]

    Grazie alla formula di eulero, $e^{ix} = \cos(x) + i \sin(x)$, è possibile esprimere un generico numero complesso $z \in \mathbb{C}$ in due forme, ovvero:
    \begin{enumerate}
        \item Forma cartesiana: $z = a + ib, \; a, b \in \mathbb{R}$
        \item Forma polare (o esponenziale): $z = \rho e^{i \theta}, \; \rho \in \mathbb{R}, \theta \in [0, 2 \pi]$
    \end{enumerate}
    E' inoltre possibile passare da una all'altra forma nei seguenti modi
    \[
        (1 \rightarrow 2) \quad z = a + ib \implies \rho = \sqrt{a^{2} + b^{2}}, \; \theta = \textnormal{atan2}\qty(\frac{b}{a})
    \]
    \[
        (2 \rightarrow 1) \quad z = \rho e^{i \theta} \implies a = \rho \cos(\theta), \; b = \rho \sin(\theta)
    \]



    \vspace{10pt}
    \subsection{Operazioni sui complessi}
        In qualità di estensione di campo, $\mathbb{C}$ è a sua volta un campo, dove addizione e moltiplicazione sono definite come segue 
        \begin{itemize}
            \item Addizione: $\forall z_{1}, z_{2} \in \mathbb{C}, \; z_{3} = z_{1} + z_{2} = (a_{1} + a_{2}) + i(b_{1} + b_{2})$ (NOTA: L'addizione è più facile da eseguire in forma cartesiana)
            \item Moltiplicazione: $\forall z_{1}, z_{2} \in \mathbb{C}, \; z_{3} = z_{1}z_{2} = (a_{1}a_{2} - b_{1}b_{2}) + i(a_{1}b_{2} + b_{1}a_{2})$ (NOTA: la moltiplicazone è più facile da eseguire in forma polare: $z_{1}z_{2} = \rho_{1}\rho_{2}e^{i(\theta_{1} + \theta_{2})}$)
        \end{itemize}
    


    \vspace{10pt}
    \subsection{Proprietà algebriche di $\mathbb{C}$}
        Sia $\mathbb{C}[z]$ l'anello dei polinomi a coefficienti complessi, è possibile allora enunciare i seguenti teoremi

        \begin{theorem}[Teorema fondamentale dell'algebra]
            Sia $P_{n}(z) \in \mathbb{C}[z]$ un polinomio a coefficienti complessi non costante.
            Questi allora ha almeno uno zero in $\mathbb{C}$, ovvero
            \[
                \exists z_{0} \in \mathbb{C} \; | \; P_{n}(z_{0}) = a_{n} (z_{0})^{n} + ... + a_{0} (z_{0})^{0} = 0
            \]
        \end{theorem}

        \begin{theorem}[Teorema di Ruffini]
            Sia $P_{n}(x) \in \mathbb{F}[x]$ un polinomio e sia $x_{0} \in \textnormal{Dom}(P_{n})$ un suo zero.
            Allora
            \[
                P_{n}(x) = \sum_{i = 0}^{n} a_{i} x^{i} = (x - x_{0}) \sum_{i = 0}^{n - 1} b_{i} x^{i} = (x - x_{0}) Q_{n - 1}(x), \quad Q_{n - 1}(x) \in \mathbb{F}[x]
            \]
        \end{theorem}

        Combinando i due teoremi, si raggiunge la conclusione

        \begin{corollary}[Corollario al Teorema fondamentale dell'algebra]
            Sia $P_{n}(z) \in \mathbb{C}[z]$ un polinomio a coefficienti complessi non costante di grado $n$.
            Allora $P_{n}(z)$ ha esattamente $n$ zeri ed è possibile esprimerlo come prodotto di polinomi di primo grado, ovvero
            \[
                P_{n}(z) = \prod_{i = 1}^{n} (z - z_{i})
            \]
        \end{corollary}

        \begin{proof}[Dimostrazione]
            Per il Teorema fondamentale dell'algebra, si ha che 
            \[
                \exists z_{1} \in \mathbb{C} \; | \; P_{n}(z_{1}) = 0
            \]

            per il Teorema di Ruffini, allora
            \[
                P_{n}(z) = (z - z_{1}) Q_{n - 1}(z)
            \]

            ma $Q_{n - 1}(z)$ soddisfa le ipotesi del Teorema fondamentale dell'algebra, il che implica che è possibile ripetere il processo iterativamente fino ad ottenere

            \[
                P_{n}(z) = \prod_{i = 1}^{n} (z - z_{i})
            \]
        \end{proof}



\newpage
\section{Morfismi complessi in dominio complesso}
    Prima ancora di definire l'idea di un morfismo complesso a dominio complesso, è necessario chiarire il concetto preliminare di \textit{dominio}
    \begin{definition}[Definizione di dominio]
        Un dominio è un qualsiasi $\Omega \subseteq \mathbb{C}$ tale per cui esso è aperto e connesso.
    \end{definition}

    Una funzione complessa a variabile complessa è una qualsiasi mappa del tipo $f: \Omega \subseteq \mathbb{C} \rightarrow \mathbb{C}$.
    Un'esempio sono alcune delle funzioni elementari, fra cui

    \begin{itemize}
        \item Esponenziale: $e^{z} = e^{a + ib} = e^{a}e^{ib} = e^{a}(\cos(b) + i \sin(b))$
        \item Logaritmo naturale: $\ln(z) = \ln(\rho e^{i \theta}) = \ln(\rho) + i\theta$
        \item Seno: $\sin(z) = \displaystyle \frac{e^{iz} - e^{-iz}}{2i}$
        \item Coseno: $\cos(z) = \displaystyle \frac{e^{iz} + e^{-iz}}{2}$
        \item Polinomi: $P_{n}(z) \in \mathbb{C}[z], \; P_{n}(z) = a_{n} z^{n} + ... + a_{0}$
    \end{itemize}

    Proiettando un morfismo complesso a variabile complessa si possono ottenre le \textit{componenti} reali di $f$, ovvero
    \[
        f(z) = u(z) + i v(z) \implies f(a + i b) = u(a, b) + iv(a, b) \quad \textnormal{con} \; u, v : \mathbb{R}^{2} \rightarrow \mathbb{R}
    \]

    \subsection{Continuità e limiti}
        Avendo la fortuna di operare in uno spazio metrico, si esprime la nozione di continuità in campo complesso nel seguente modo:

        \begin{definition}[Definizione di continuità in un punto]
            Sia $f: \Omega \subseteq \mathbb{C} \rightarrow \mathbb{C}$ una funzione complessa a variabile complessa e sia $z_{0} \in \Omega$ un punto del dominio.
            $f$ si dice continua in $z_{0}$ se e solo se per ogni $\epsilon > 0$ esiste un $\delta > 0$ tale per cui $f(B_{\delta}(z_{0}) \cap \Omega) \subseteq B_{\epsilon}(f(z_{0}))$.
        \end{definition}

        Dalla precedente definizione definiamo la continuità (nel proprio dominio) di una funzione come continuità in tutti i punti del suo dominio.
        Inoltre, non è difficile mostrare che se la funzione $f$ è continua in un punto, anche le sue proiezioni $u, v$ lo sono in quel punto.
        Infatti
        \[
            |f(z) - f(z_{0})| < \epsilon \implies (u(z) - u(z_{0}))^{2} + (v(z) - v(z_{0}))^{2} < \epsilon^{2} \implies |u(z) - u(z_{0})| < \sqrt{\epsilon^{2} - (v(z) - v(z_{0}))^{2}}
        \]
        \[
            |f(z) - f(z_{0})| < \epsilon \implies (u(z) - u(z_{0}))^{2} + (v(z) - v(z_{0}))^{2} < \epsilon^{2} \implies |v(z) - v(z_{0})| < \sqrt{\epsilon^{2} - (u(z) - u(z_{0}))^{2}}
        \]
    
        Sulla base delle precedenti considerazioni definiamo anche il concetto di punto di accumulazione e di limite.

        \begin{definition}[Punto di accumulazione]
            Dato un insieme $D \subseteq \mathbb{C}$, un punto $z_{0} \in \mathbb{C}$ si dice punto di accumulazione di $D$ se
            \[
                \forall \epsilon > 0, \; B_{\epsilon}(z_{0}) \setminus \{z_{0}\} \cap D \neq \emptyset
            \]
        \end{definition}

        \begin{definition}[Limite]
            Data una funzione $f: \Omega \subseteq \mathbb{C} \rightarrow \mathbb{C}$, si dice che $L \in \mathbb{C}$, punto di accumulazione dell'immagine, è limite di $f$ che tende a $z_{0} \in \mathbb{C}$, punto di accumulazione del dominio, qualora 
            \[
                \forall \epsilon > 0, \; \exists \delta > 0 \; \textnormal{tale che} \; f(B_{\delta}(z_{0}) \setminus \{z_{0}\} \cap \Omega) \subseteq B_{\epsilon}(L)
            \]
        \end{definition}

        La principale differenza fra la definizione di limite e la definizione di continuità in un punto è che nel caso del limite la funzione non deve necessariamente \textit{raggiungere} $L$, ovvero $f(z) = L$ per qualche $z \in \Omega$.
        E' possibile allora riscrivere la definizione di continuità attraverso la definizione di limite.
        \[
            f \; \textnormal{è continua in} \; z_{0} \iff \lim_{z \rightarrow z_{0}} f(z) = f(z_{0})
        \]

        Un'altra, particolarmente importante, osservazione è che, se le componenti di una funzione sono continue in un punto, allora anche la funzione è continua in quel punto.
        Infatti, per ogni $\epsilon > 0$ si ha che esistono $\delta_{1}, \delta_{2} > 0$, per $u$ e $v$ rispettivamente, tali per cui i punti nel dominio entro $\min(\delta_{1}, \delta_{2})$ da $z_{0}$ cadono interamente entro $\sqrt{2} \epsilon$ da $f(z_{0})$


    \subsection{Calcolo differenziale}
        \begin{definition}[Derivata di una fuzione complessa]
            Sia $f: \Omega \subseteq \mathbb{C} \rightarrow \mathbb{C}$ con $\Omega$ aperto e $z_{0} \in \Omega$ un generico punto del dominio.
            Si dice derivata di $f$ in $z_{0}$ la seguente espressione
            \[
                f'(z_{0}) = \lim_{z \rightarrow z_{0}} \frac{f(z) - f(z_{0})}{z - z_{0}} = \lim_{z \rightarrow z_{0}} \frac{u(z) - u(z_{0})}{z - z_{0}} + i \lim_{z \rightarrow z_{0}} \frac{v(z) - v(z_{0})}{z - z_{0}} = u'(z_{0}) + i v'(z_{0})
            \]
        \end{definition}

        \begin{definition}[Funzione olomorfa]
            Sia $f: \Omega \subseteq \mathbb{C} \rightarrow \mathbb{C}$ con $\Omega$ aperto.
            $f$ si dice olomorfa se questa ammette derivata su tutto $\Omega$.
        \end{definition}

        \begin{definition}[Funzione intera]
            Sia $f: \mathbb{C} \rightarrow \mathbb{C}$.
            $f$ si dice intera se questa è olomorfa.
        \end{definition}

        In generale, possiamo trattare funzioni complesse a variabile complessa come speciali campi vettoriali da e verso $\mathbb{R}^{2}$, e di conseguenza si applicano le stesse definizioni di differenziabilità, ovvero che $u, v$ devono essere differenziabili.
        Ricordiamo che condizione sufficiente per la differenziabilità di un campo scalare è che le derivate parziali esistano e siano continue.

        \vspace{10pt}
        \subsubsection{Equazioni di Cauchy-Riemann}
            Inoltre, una importante conseguenza del fatto che la definizione di limite deve valere lungo qualsiasi direzione si approccia $z_{0}$ è che le seguenti condizioni devono essere vere contemporaneamente
            \begin{enumerate}
                \item Approcciando $z_{0}$ lungo l'asse reale, abbiamo che 
                \[
                    f'(z_{0}) = \lim_{z \rightarrow z_{0}} \frac{f(z) - f(z_{0})}{z - z_{0}} = \lim_{\Delta x \rightarrow 0} \frac{f(x_{0} + \Delta x, y_{0}) - f(x_{0}, y_{0})}{\Delta x} 
                \]
                Che espandendo diventa
                \[
                    \lim_{\Delta x \rightarrow 0} \frac{u(x_{0} + \Delta x, y_{0}) - u(x_{0}, y_{0})}{\Delta x} + \; i \lim_{\Delta x \rightarrow 0} \frac{v(x_{0} + \Delta x, y_{0}) - v(x_{0}, y_{0})}{\Delta x} = \frac{\partial u}{\partial x} + i \frac{\partial v}{\partial x}
                \]

                \item Approcciando $z_{0}$ lungo l'asse immaginario, abbiamo che
                \[
                    f'(z_{0}) = \lim_{z \rightarrow z_{0}} \frac{f(z) - f(z_{0})}{z - z_{0}} \lim_{\Delta y \rightarrow 0} \frac{f(x_{0}, y_{0} + \Delta y) - f(x_{0}, y_{0})}{i \Delta y} 
                \]
                Che espandendo diventa
                \[
                    \lim_{\Delta y \rightarrow 0} \frac{u(x_{0}, y_{0} + \Delta y) - u(x_{0}, y_{0})}{i \Delta y} + \; i \lim_{\Delta y \rightarrow 0} \frac{v(x_{0}, y_{0} + \Delta y) - v(x_{0}, y_{0})}{i \Delta y} = \frac{\partial v}{\partial y} - i \frac{\partial u}{\partial y}
                \]
            \end{enumerate}

            Dato che queste due condizioni devono essere vere simultaneamente, abbiamo che 
            \[
                \begin{cases}
                    \displaystyle \frac{\partial u}{\partial x} = \frac{\partial v}{\partial y} \\
                    \\
                    \displaystyle \frac{\partial v}{\partial x} = -\frac{\partial u}{\partial y} \\
                \end{cases} 
            \]

            \begin{theorem}[Teorema sulla differenziabilità di funzioni complesse]
                Sia $f: \Omega \subseteq \mathbb{C} \rightarrow \mathbb{C}$ una funzione con $\Omega$ aperto.
                Allora se seguenti condizioni sono equivalenti
                \begin{enumerate}
                    \item $f$ è differenziabile in un punto $z_{0}$
                    \item Le componenti $u, v$ di $f$ sono differenziabili in $z_{0}$, e valgono le condizioni di Cauchy-Riemann nello stesso
                \end{enumerate}
            \end{theorem}

            \begin{proposition}
                Sia $f: \Omega \subseteq \mathbb{C} \rightarrow \mathbb{C}$ una funzione olomorfa.
                Se $u, v$ sono di classe $C^{2}(\Omega)$, allora queste soddisfano l'equazione di Laplace.
            \end{proposition}

            \begin{proof}[Dimostrazione]
                Da Cauchy-Riemann, abbiamo che
                \[
                    \frac{\partial^{2} u}{\partial x^{2}} = \frac{\partial^{2} v}{\partial y \partial x} 
                \]

                Ricordando il Teorema di Schwarz, il quale postula che, per campi scalari di classe $C^{2}$
                \[
                    \frac{\partial^{2} \{u, v\}}{\partial x \partial y} = \frac{\partial}{\partial y} \qty(\frac{\partial \{u, v\}}{\partial x}) = \frac{\partial}{\partial x} \qty(\frac{\partial \{u, v\}}{\partial y})
                \]

                Di conseguenza si ha che
                \[
                    \frac{\partial^{2} u}{\partial x^{2}} = \frac{\partial^{2} v}{\partial y \partial x} = \frac{\partial}{\partial y} \qty(\frac{\partial v}{\partial x}) = -\frac{\partial^{2} u}{\partial y^{2}} \implies \frac{\partial^{2} u}{\partial y^{2}} + \frac{\partial^{2} u}{\partial x^{2}} = 0 \
                \]
            \end{proof}

            \begin{proposition}
                Sia $u: \mathbb{R}^{2} \rightarrow \mathbb{R}$ una funzione armonica (che soddisfa l'eq. di Laplace).
                Allora esiste una classe di funzioni $v: \mathbb{R}^{2} \rightarrow \mathbb{R}$ armoniche tali per cui $f = u + i v$ è intera, chiamate complementi armonici di $u$.
            \end{proposition}
        


    \subsection{Curve ed integrali}
        \begin{definition}[Definizione di curva]
            Una mappa $\gamma: [a, b] \rightarrow \mathbb{C}$ è detta curva se è continua e di classe $C^{1}$ a tratti su $[a, b]$.
            $S(\gamma) = \gamma([a, b])$ è detto supporto di $\gamma$.
        \end{definition}

        \begin{definition}[Definizione di curva regolare]
            Una curva $\gamma: [a, b] \rightarrow \mathbb{C}$ è detta regolare se $\forall t \in [a, b]$ abbiamo che $\gamma'(t) \neq 0$.
        \end{definition}

        \begin{definition}[Definizione di curva semplice]
            Una curva $\gamma: [a, b] \rightarrow \mathbb{C}$ è detta semplice se è iniettiva, ovvero
            \[
                \forall t_{1}, t_{2} \in [a, b], \; t_{1} \neq t_{2} \implies \gamma(t_{1}) \neq \gamma(t_{2})
            \]
        \end{definition}

        \begin{definition}[Definizione di curva chiusa]
            Una curva $\gamma: [a, b] \rightarrow \mathbb{C}$ è detta chiusa se $\gamma(a) = \gamma(b)$.
        \end{definition}

        Un esempio comune di curve sono le curve circonferenziali, ovvero $\gamma: [0, 2\pi] \rightarrow \mathbb{C}$ con $\gamma(t) = z_{0} + re^{it}$.

        \subsubsection{Operazioni sulle curve}
            \begin{enumerate}
                \item Concatenazione. Siano $\gamma_{1}: [a, b] \rightarrow \mathbb{C}$, $\gamma_{2}: [b, c] \rightarrow \mathbb{C}$ due curve tali per cui $\gamma_{1}(b) = \gamma_{2}(b)$, possiamo allora definire una nuova curva, $(\gamma_{1} * \gamma_{2}): [a, c] \rightarrow \mathbb{C}$ tale che
                \[
                    (\gamma_{1} * \gamma_{2})(t) := 
                    \begin{cases} 
                        \gamma_{1}(t), & t \in [a, b] \\
                        \gamma_{2}(t), & t \in (b, c]
                    \end{cases}
                \]

                \item Riparametrizzazione. Sia $\gamma: [a, b] \rightarrow \mathbb{C}$ una curva e sia $\phi: [c, d] \rightarrow [a, b]$ un orologio di classe $C^{1}$ strettamente monotono.
                Allora $\gamma \circ \phi: [c, d] \rightarrow \mathbb{C}$ è una curva il cui supporto è identico alla curva $\gamma$ ma il cui verso di percorrenza è dettato dalla monotonia dell'orologio: se questi è strettamente crescente, allora il verso è lo stesso di $\gamma$; se questi è strettamente decrescente, allora il verso è opposto a quello di $\gamma$.
            \end{enumerate}
            
        \subsubsection{Curve di Jordan}
            \begin{definition}[Definizione di curva di Jordan]
                Una curva $\gamma: [a, b] \rightarrow \mathbb{C}$ chiusa è detta di Jordan se questa è semplice, fatta eccezione per gli estremi del dominio. 
            \end{definition}

            Le curve di Jordan sono particolarmente importanti in quanto su di esse è basato un importante teorema sulla ripartizione topologica dei piani, fra cui figura quello \textit{complesso}.

            \begin{theorem}[Teorema delle curve (complesse) di Jordan]
                Sia $\gamma: [a, b] \rightarrow \mathbb{C}$ una curva di Jordan.
                Allora $\mathbb{C} \setminus S(\gamma)$ è un insieme aperto esprimibile come unione disgiunta di due domini: uno limitato, l'interno di $\gamma$ $(\textnormal{Int}(\gamma))$, ed uno illimitato, l'esterno di $\gamma$ $(\textnormal{Ext}(\gamma))$.
            \end{theorem}

            Grazie alla nozione di curve di Jordan, ci è permesso di esprimere un concetto di vitale importanza, la proprieta di \textit{connessione semplice}.

            \begin{definition}[Definizione di dominio semplicemente connesso]
                Sia $D \subseteq \mathbb{C}$ un dominio.
                Questi è detto semplicemente connesso se e solo se tutte le curve di Jordan $\gamma$ tali per cui $S(\gamma) \subseteq D$ hanno il loro interno in $D$, ovvero
                \[
                    \textnormal{Int}(\gamma) \subseteq D
                \]
            \end{definition}

            \begin{definition}[Definizione di dominio con bordo]
                Sia $D \subseteq \mathbb{C}$ un dominio, e sia $\partial D$ la sua frontiera. 
                $D$ si dice allora un dominio dotato di bordo se la sua frontiera è data da un'unione disgiunta di supporti di curve di Jordan, ovvero
                \[
                    \partial D = \bigsqcup_{i} S(\gamma_{i})
                \]

                NB: Il bordo è detto orientato positivamente ($\partial D^{+}$) se, percorrendo ogni singola $\gamma_{i}$, si lasciano i punti del dominio "alla propria sinistra".
            \end{definition}

        \subsubsection{Integrali di linea}
            Data una curva $\gamma: [a, b] \rightarrow \mathbb{C}$, $\gamma(t) = x(t) + i y(t)$ si definisce allora l'integrale di $\gamma: [a, b] \rightarrow \mathbb{C}$ come
            \[
                \int_{a}^{b} \gamma \; dt := \int_{a}^{b} x(t) \; dt + i \int_{a}^{b} y(t) \; dt
            \]

            Data anche una funzione $f: D \subseteq \mathbb{C} \rightarrow \mathbb{C}$ definita su un dominio che contiene il supporto di $\gamma$, è possibile definire il seguente tipo di integrale
            \[
                \int_{\gamma} f \: dz := \int_{a}^{b} f(\gamma) \: \gamma' \:  dt = \int_{a}^{b} \qty[u(x, y) + i v(x, y)](x' + i y') dt = 
            \]
            \[
                \int_{a}^{b} \qty[u(x, y)x' - v(x, y)y'] \: dt + i \int_{a}^{b} \qty[v(x, y)x' + u(x, y)y'] \: dt = \int_{\gamma} (u, -v) \cdot d\gamma + i \int_{\gamma} (v, u) \cdot d\gamma
            \]
            dove $d\gamma = (x' dt, y' dt)$

            \vspace{10pt}
            Questo tipo di integrale condivide molte proprietà degli integrali di linea su $\mathbb{R}^{n}$, ovvero
            \begin{enumerate}
                \item Linearità:
                \[
                    \int_{\gamma} (f + g) \: dz = \int_{\gamma} f \: dz + \int_{\gamma} g \: dz
                \]

                \item Additività:
                \[
                    \int_{(\gamma_{1} * \gamma_{2})} f \: dz = \int_{\gamma_{1}} f \: dz + \int_{\gamma_{2}} f \: dz
                \]

                \item Riparametrizzazione: Sia $\phi: [c, d] \rightarrow [a, b]$ un orologio di classe $C^{1}$ strettamente crescente, allora
                \[
                    \int_{\gamma \circ \phi} f \: dz = \int_{\gamma} f \: dz
                \]
                Qualora $\phi$ fosse strettamente decrescente, si otterrebbe
                \[
                    \int_{\gamma \circ \phi} f \: dz = - \int_{\gamma} f \: dz
                \]
            \end{enumerate}

            Sugli integrali di linea è impostato uno dei teoremi più importanti dell'analisi complessa: il Teorema di Cauchy-Goursat

            \vspace{10pt}
            \begin{theorem}[Teorema di Cauchy-Goursat]
                Sia $f: D \subseteq \mathbb{C} \rightarrow \mathbb{C}$ una funzione olomorfa definita su un dominio $D$ semplicemente connesso e sia $\gamma: [a, b] \rightarrow D$ una curva di Jordan.
                Allora
                \[
                    \oint_{\gamma} f \: dz = 0
                \]
            \end{theorem}

            \begin{proof}[Dimostrazione]
                Si prenda in considerazione la forma generale dell'integrale
                \[
                    \oint_{\gamma} f \: dz = \oint_{\gamma} (u, -v) \cdot d\gamma + i \oint_{\gamma} (v, u) \cdot d\gamma
                \]
                    
                Si osservi allora il comportamento della Jacobiana dei seguenti campi: $(u, -v)$ e $(v, u)$
                \[
                    J((u, -v)) =
                    \begin{bmatrix}
                        \frac{\partial u}{\partial x} & \frac{\partial u}{\partial y} \\[1em]
                        \frac{\partial (-v)}{\partial x} & \frac{\partial (-v)}{\partial y} \\
                    \end{bmatrix} 
                \]
                \[
                    J((v, u)) = 
                    \begin{bmatrix}
                        \frac{\partial v}{\partial x} & \frac{\partial v}{\partial y} \\[1em]
                        \frac{\partial u}{\partial x} & \frac{\partial u}{\partial y} \\
                    \end{bmatrix}
                \]

                Data la natura olomorfa di $f$, possiamo sostituire in accordo con le equazioni di Cauchy-Riemann
                \[
                    J((u, -v)) =
                    \begin{bmatrix}
                        \frac{\partial u}{\partial x} & -\frac{\partial v}{\partial x} \\[1em]
                        \frac{\partial (-v)}{\partial x} & \frac{\partial (-v)}{\partial y} \\
                    \end{bmatrix} 
                \]
                \[
                    J((v, u)) = 
                    \begin{bmatrix}
                        \frac{\partial v}{\partial x} & \frac{\partial u}{\partial x} \\[1em]
                        \frac{\partial u}{\partial x} & \frac{\partial u}{\partial y} \\
                    \end{bmatrix}
                \]

                Abbiamo appurato allora che le Jacobiane di $(u, -v)$ e $(v, u)$ sono simmetriche, il che implica che i corrispettivi campi sono localmente conservativi.
                Tale considerazione, congiunta con il fatto che $(u, -v)$ e $(v, u)$ sono definiti su un dominio semplicemente connesso, permettono di asserire che
                \[
                    \oint_{\gamma} (u, -v) \cdot d\gamma = 0 \; \textnormal{e} \; \oint_{\gamma} (v, u) \cdot d\gamma = 0 \implies \oint_{\gamma} f \: dz = 0
                \]
            \end{proof}

            \begin{corollary}[Corollario al Teorema di Cauchy-Goursat]
                Sia $f: D \subseteq \mathbb{C} \rightarrow \mathbb{C}$ una funzione olomorfa definita su un dominio $D$ semplicemente connesso.
                Siano $\gamma_{1}: [a, b] \rightarrow D$, $\gamma_{2}: [a, b] \rightarrow D$ due curve semplici tali per cui $\gamma_{1}(a) = \gamma_{2}(a)$ e $\gamma_{1}(b) = \gamma_{2}(b)$ e per ogni $t \in (a, b), \; \gamma_{1}(t) \neq \gamma_{2}(t)$.
                Allora
                \[
                    \int_{\gamma_{1}} f \: dz = \int_{\gamma_{2}} f \: dz
                \]
            \end{corollary}

            \begin{proof}[Dimostrazione]
                Si consideri la concatenazione $(-\gamma_{1} * \gamma_{2})$. 
                Questa è una curva di Jordan.
                Per il Teorema di Cauchy-Goursat
                \[
                    \oint_{(-\gamma_{1} * \gamma_{2})} f \: dz = - \int_{\gamma_{1}} f \: dz + \int_{\gamma_{2}} f \: dz = 0 \implies \int_{\gamma_{1}} f \: dz = \int_{\gamma_{2}} f \: dz
                \]
            \end{proof}

            \begin{theorem}[Teorema di Cauchy-Goursat generalizzato]
                Sia $f: \Omega \subseteq \mathbb{C} \rightarrow \mathbb{C}$ una funzione olomorfa.
                Sia $D \subseteq \Omega$ un dominio il cui bordo orientato positivamente è denotato da $\partial D$.
                Allora
                \[
                    \oint_{\partial D} f \: dz = \sum_{i = 1}^{n} \oint_{\gamma_{i}} f \: dz = 0
                \]
            \end{theorem}

            \begin{corollary}[Corollario al Teorema di Cauchy-Goursat generalizzato]
                Sia $f: D \setminus \{z_{1}, ..., z_{n}\} \subseteq \mathbb{C} \rightarrow \mathbb{C}$ una funzione olomorfa.
                Siano $\gamma_{1}, \gamma_{2}$ due curve di Jordan che racchiudano entrambe un qualsiasi $B \in \mathcal{P}(\{z_{i}\}_{i = 1}^{n})$, con la condizione che $S(\gamma_{\{1, 2\}}) \cap B = \emptyset$.
                Allora
                \[
                    \int_{\gamma_{1}} f \: dz = \int_{\gamma_{2}} f \: dz
                \]
            \end{corollary}

        

    \vspace{10pt}
    \subsection{Analiticità e residui di funzioni}

        \begin{definition}[Formula di Cauchy]
            Sia $f: D \subseteq \mathbb{C} \rightarrow \mathbb{C}$ una funzione olomorfa definita su un dominio semplicemente connesso, e sia $\gamma: [a, b] \rightarrow D$ una curva di Jordan orientata positivamente.
            Dato un $z_{0} \in D$, si ha che
            \[
                \oint_{\gamma} \frac{f}{(z - z_{0})} \: dz = 2 \pi i f(z_{0})
            \]
        \end{definition}

        \begin{proof}[Dimostrazione]
            Si consideri la seguente curva $\gamma(t) = z_{0} + re^{it}$.
            L'integrale diventa allora
            \[
                \oint_{\gamma} \frac{f}{(z - z_{0})} \: dz = \oint_{\gamma} \frac{f(z) - f(z_{0}) + f(z_{0})}{(z - z_{0})} \: dz = f(z_{0}) \oint_{\gamma} \frac{1}{z - z_{0}} \: dz + \oint_{\gamma} \frac{f(z) - f(z_{0})}{(z - z_{0})} \: dz
            \]

            La prima parte, sappiamo valere $f(z_{0}) \cdot 2 \pi i$.
            Per quanto riguarda la seconda parte
            \[
                \oint_{\gamma} \frac{f(z) - f(z_{0})}{(z - z_{0})} \: dz = \int_{0}^{2 \pi} \frac{f(z_{0} + re^{it}) - f(z_{0})}{re^{it}} ire^{it} \: dt = i \int_{0}^{2 \pi} f(z_{0} + re^{it}) - f(z_{0}) \: dt
            \]

            Per il Lemma di stima, abbiamo che
            \[
                0 \leq \abs{\int_{0}^{2 \pi} f(z) - f(z_{0}) \: dt} \leq \max_{z \in S(\gamma)}(\abs{f(z) - f(z_{0})}) \cdot 2 \pi
            \]

            Per il Corollario al Teorema di Cauchy-Goursat generalizzato, la soprastante disuguaglianza deve essere verificata per qualsiasi $r \in \mathbb{R}^{+}$ e che $f$, in qualità di funzione olomorfa, è anche continua, si osservi cosa accade per $r \rightarrow 0$
            \[
                \lim_{r \rightarrow 0} \max_{z \in S(\gamma)}(\abs{f(z) - f(z_{0})}) = 0 \implies 0 \leq \abs{\int_{0}^{2 \pi} f(z_{0} + re^{it}) - f(z_{0}) \: dt} \leq 0
            \]
            \[
                \oint_{\gamma} \frac{f}{(z - z_{0})} \: dz = f(z_{0}) \oint_{\gamma} \frac{1}{z - z_{0}} \: dz + \oint_{\gamma} \frac{f(z) - f(z_{0})}{(z - z_{0})} \: dz = 2 \pi f(z_{0}) + 0
            \]
        \end{proof}
                
        \begin{definition}[Formula di Cauchy generalizzata]
            Sia $f: D \subseteq \mathbb{C} \rightarrow \mathbb{C}$ una funzione olomorfa definita su un dominio semplicemente connesso, e sia $\gamma: [a, b] \rightarrow D$ una curva di Jordan orientata positivamente.
            Dato un $z_{0} \in D$, si ha che
            \[
                \oint_{\gamma} \frac{f}{(z - z_{0})^{n + 1}} \: dz = \frac{2 \pi i}{n!} f^{(n)}(z_{0})
            \]
        \end{definition}

        \vspace{10pt}
        Armati di questa collezione di importanti teoremi e corollari, possiamo enunciare la seguente formula

        \begin{theorem}[Teorema di Liouville]
            Sia $f: \mathbb{C} \rightarrow \mathbb{C}$ una funzione intera e limitata.
            Allora $f$ è costante.
        \end{theorem}

        \begin{proof}[Dimostrazione]
            Si selezioni un generico $z_{0} \in \mathbb{C}$ e una circonferenza di raggio $r \in \mathbb{R}^{+}$. Per la formula generalizzata di Cauchy, si ha che
            \[
                \int_{0}^{2\pi} \frac{f(z_{0} + re^{it})}{((z_{0} + re^{it}) - z_{0})^{2}} ire^{it} \: dt = 2 \pi i f'(z_{0}) \implies \frac{1}{2 \pi r} \int_{0}^{2 \pi} f(z_{0} + re^{it}) e^{-it} dt = f'(z_{0})
            \]

            Allora
            \[
                \abs{f'(z_{0})} = \frac{1}{2 \pi r} \abs{\int_{0}^{2 \pi} f(z_{0} + re^{it}) e^{-it} dt} \leq \frac{1}{2 \pi r} \int_{0}^{2 \pi} \abs{f(z_{0} + re^{it}) e^{-it}} dt \leq \frac{k}{2 \pi r}
            \]

            Dato che la seguente espressione
            \[
                \abs{f'(z_{0})} \leq \frac{k}{2 \pi r}
            \]
            deve essere valida per qualsiasi $r \in \mathbb{R}^{+}$ e per qualsiasi $z_{0} \in \mathbb{C}$, si ha che
            \[
                \abs{f'(z_{0})} = 0, \; \forall z_{0} \in \mathbb{C}
            \]
        \end{proof}

        \subsubsection{Sviluppi di Laurent}
            \begin{theorem}[Teorema di Laurent]
                Sia $f: B_{r}(z_{0}) \setminus \{z_{0}\} \rightarrow \mathbb{C}$ una funzione olomorfa.
                Allora
                \[
                    f(z) = \sum_{-\infty}^{+\infty} a_{k} (z - z_{0})^{k}
                \]
                dove 
                \[
                    a_{k} = \frac{1}{2 \pi i} \oint_{\partial B_{\delta}(z_{0})} \frac{f}{(z - z_{0})^{k + 1}} \: dz \quad \textnormal{con} \; \delta < r
                \]
            \end{theorem}

            \begin{definition}[Definizione di residuo]
                Sia $f: B_{r}(z_{0}) \setminus \{z_{0}\} \rightarrow \mathbb{C}$ una funzione olomorfa.
                Si dice allora il coefficiente 
                \[
                    a_{-1} = \frac{1}{2 \pi i} \oint_{\partial B_{\delta}(z_{0})} f(z) \: dz \quad \textnormal{con} \; \delta < r
                \]
                residuo di $f$ in $z_{0}$ e si denota nel seguente modo
                \[
                    a_{-1} = \textnormal{Res}(f, z_{0}) \implies \oint_{\partial B_{\delta}(z_{0})} f(z) \: dz = 2 \pi i \: \textnormal{Res}(f, z_{0})
                \]
            \end{definition}


            \vspace{10pt}
            La serie di Laurent si dimostra un'essenziale strumento nell'analisi complessa.
            La parte della serie avente $k \geq 0$ si dice parte regolare, mentre la parte determinata da $k < 0$ si dice parte principale.
            La serie di Laurent ci permette di caratterizzare le singolarità di una funzione:
            \begin{itemize}
                \item Se non vi è presente parte principale, allora $f = \sum_{k = 0}^{\infty} a_{k} (z - z_{0})^{k}$, che è una semplice serie di Taylor.
                E' allora possibile determinare il valore di $f$ in $z_{0}$, in quanto $f(z_{0}) = a_{0}$, e la singolarità è definita eliminabile.

                \item Se vi sono presenti $m$ termini principali, allora $f = \sum_{k = 1}^{m} a_{-k} (z - z_{0})^{-k} + \sum_{k = 0}^{\infty} a_{k} (z - z_{0})^{k}$, e $z_{0}$ è definito polo di ordine $m$.

                \item Se vi sono presenti infiniti termini principali, allora $f = \sum_{k = -\infty}^{\infty} a_{k}(z - z_{0})^{k}$, e $z_{0}$ è detta singolarità essenziale.
            \end{itemize}


            \begin{theorem}[Teorema dei residui]
                Sia $f: D \setminus \{z_{1}, ..., z_{n}\} \subseteq \mathbb{C} \rightarrow \mathbb{C}$ una funzione olomorfa e sia $D$ un dominio semplicemente connesso.
                Sia $\gamma: [a, b] \rightarrow D$ una curva di Jordan che racchiuda un numero finito di singolarità $\{z_{1}, ..., z_{p}\} \subseteq \textnormal{Int}(\gamma)$, $p \leq n$, in $\textnormal{Int}(\gamma)$.
                Allora
                \[
                    \oint_{\gamma} f \: dz = 2\pi i \sum_{k = 1}^{p} \textnormal{Res}(f, z_{k})
                \]
            \end{theorem}

            \begin{proof}[Dimostrazione]
                Si immagini di costruire un nuovo dominio, $U \subseteq D$, tale per cui il bordo $\partial U$ orientato positivamente è determinato da $\gamma$ e, per ogni singolarità $z_{k} \in $ una circonferenza $\gamma_{k}(t) = z_{k} + r_{k}e^{-it}$ con $r_{k}$ selezionato in modo da $S(\gamma_{k}) \subseteq D$.
                Per il Teorema di Cauchy-Goursat generalizzato, si ottiene
                \[
                    \oint_{\partial U} f \: dz = \oint_{\gamma} f \: dz + \sum_{k = 1}^{p} \oint_{\gamma_{k}} f \: dz = 0 \implies \oint_{\gamma} f \: dz = - \sum_{k = 1}^{p} \oint_{\gamma_{k}} f \: dz
                \]

                Per ogni $\gamma_{k}$ vale la seguente uguaglianza, in quanto orientate in senso orario
                \[
                    \oint_{\gamma_{k}} f \: dz = - 2 \pi i \: \textnormal{Res}(f, z_{k})
                \]

                Allora
                \[
                    \oint_{\gamma} f \: dz = 2 \pi i \sum_{k = 1}^{p} \textnormal{Res}(f, z_{k})
                \]
            \end{proof}
            

    \newpage
    \subsection{Valore principale di Cauchy}
        \begin{definition}[Definizione di valore principale]
            Sia $f: \mathbb{C} \rightarrow \mathbb{C}$ una funzione complessa a valori complessi e sia $z_{0}$ un polo di $f$. 
            Sia $\gamma: [a, b] \rightarrow \mathbb{C}$ una curva regolare tale per cui, per un valore $a < c < b$, si ha che $\gamma(c) = z_{0}$.
            Si definisce allora valore principale dell'integrale
            \[
                \int_{\gamma} f(z) \: dz
            \]
            la seguente espressione
            \[
                \pv{\int_{\gamma} f(z) \: dz} = \lim_{\epsilon \rightarrow 0^{+}} \qty( \int_{a}^{c - \epsilon} f(\gamma(t)) \gamma'(t) \: dt + \int_{c + \epsilon}^{b} f(\gamma(t)) \gamma'(t) \: dt)
            \]
        \end{definition}


\end{document}